\section{Conclusions}\label{conclusions}

This study represents a first attempt to integrate RL with DBR for
production scheduling. The RL model was developed using TOC concepts and
trained within a DES environment modeled on DBR studies, simulating a
job shop with stochastic processes. Modifications, including machine
setups and equipment breakdowns, were introduced to evaluate the
solution's performance in scenarios with high variability. The results
showed that RL achieved the same level of productivity as DBR but with
lower inventory levels.

The primary contribution of this research lies in enhancing production
scheduling efficiency by integrating RL with TOC principles and
decomposing the SB into WIP and FG. Additionally, the analysis of RL
behavior identified WIP as the key variable influencing product release
decisions, highlighting opportunities to refine DBR by incorporating WIP
more prominently. While the exact influence of WIP on RL actions could
not be fully determined, integrating RL with DBR significantly improved
the production system by reducing inventory and shortening throughput
times. Furthermore, incorporating TOC concepts into RL modeling enabled
effective decision-making without the need for exhaustive production
monitoring.

Regarding limitations, this study's simulation model is theoretical,
incorporating simplified representations of reality that exclude
operational costs and system gains (Goldratt, 2006a). Consequently,
economic analyses, including the TOC-proposed prioritization of product
mix, were not performed. Additionally, the MTS operating mode used in
this study did not account for the potential to operate in an MTO
framework. Future research opportunities include implementing the
proposed solution in more complex environments using empirical data,
comparing the RL model with alternative approaches, and further
exploring DBR by isolating WIP and FG as individual decision-making
variables.

This study adopted a mechanism-based perspective to explore whether and
how TOC can enhance Lean production theory in the context of optimizing
a welding line in a manufacturing company. We applied Lean and TOC
principles to design and implement an intervention that increases the
economic efficiency of this line by more than 230 percent. The
intervention challenged three assumptions underlying Lean implementation
and sustainment: (i) inventory and capacity buffers are cost equivalent;
(ii) reducing variability is the primary way to synchronize production
with demand; and (iii) continuous improvement inherently leads to
efficiency improvements. It also challenges the TOC assumption that
increased demand has no limit, but reduced inventory and operating
expenses do. As such, our study sheds light on the rationale for
economic efficiency-centered design that goes beyond the Lean and TOC
lenses. This rationale is grounded in the following kinematic
mechanisms: (i) measuring production systems in terms of economic
efficiency; (ii) synchronizing production and demand by eliminating
obvious waste and reducing and/or buffering variability; and (iii)
managing progress toward the economic efficiency frontier by means of
system constraints.
